% =====================================================================
% House-style preamble for the SUBMISSION PDF.
% White background (journal requirement) but house typography and layout:
% Georgia display headings + Arial body, deep-teal headings, terracotta accent
% for links, generous margins. Rendered with tectonic/XeLaTeX so the fonts
% resolve via fontspec.
%
% Fonts are chosen for REPRODUCIBILITY: the house faces are Georgia (display)
% and Arial (body), but those are proprietary and absent on CI runners, so we
% prefer the open, metric-compatible clones Gelasio (= Georgia) and Arimo
% (= Arial), falling back to the real Microsoft fonts when present locally.
% Both builds are visually 1:1 because the clones share the originals' metrics.
% =====================================================================

% --- Brand colours (white page, so background stays default white) ---
\usepackage{xcolor}
\definecolor{houseTeal}{HTML}{163139}
\definecolor{houseTerracotta}{HTML}{C16A3C}
\definecolor{houseRule}{HTML}{D9D2C4}

% --- Fonts: Arimo/Arial for body, Gelasio/Georgia for display headings ---
% Prefer the real house fonts when installed (local macOS); otherwise use the
% metric-compatible open clones that CI installs. Times/Helvetica are a last
% resort so a missing-font build still produces output rather than erroring.
%
% Georgia/Arial have proper OS-level style linking, so the bare family name is
% enough. The Gelasio clone is built as four static files by
% scripts/build_house_fonts.py; automatic bold-matching by family name is
% unreliable across CoreText/fontconfig, so its four faces are named EXPLICITLY.
\usepackage{fontspec}
\IfFontExistsTF{Arial}{\setmainfont{Arial}}{%
  \IfFontExistsTF{Arimo}{\setmainfont{Arimo}}{\setmainfont{Helvetica Neue}}}
% If the Gelasio statics exist in ./fonts (built by build_house_fonts.py, as CI
% does), reference them BY PATH — name-based bold matching is unreliable across
% CoreText/fontconfig, but an explicit Path= resolves every face deterministically.
\IfFileExists{./fonts/Gelasio-Bold.ttf}{%
  \newfontfamily\houseDisplay{Gelasio}[%
    Path=./fonts/, Extension=.ttf,%
    UprightFont={Gelasio-Regular}, BoldFont={Gelasio-Bold},%
    ItalicFont={Gelasio-Italic}, BoldItalicFont={Gelasio-BoldItalic}]%
}{%
  % No clone present (e.g. local build): use real Georgia, else Times.
  \IfFontExistsTF{Georgia}{\newfontfamily\houseDisplay{Georgia}}%
    {\newfontfamily\houseDisplay{Times New Roman}}%
}

% --- Headings in Georgia, deep teal, sentence case, left-aligned ---
\usepackage{titlesec}
\titleformat{\section}
  {\houseDisplay\bfseries\LARGE\color{houseTeal}}{\thesection}{0.6em}{}
\titleformat{\subsection}
  {\houseDisplay\bfseries\Large\color{houseTeal}}{\thesubsection}{0.5em}{}
\titleformat{\subsubsection}
  {\houseDisplay\bfseries\large\color{houseTeal}}{\thesubsubsection}{0.5em}{}
\titlespacing*{\section}{0pt}{2.2ex plus 1ex minus .2ex}{1.1ex plus .2ex}
\titlespacing*{\subsection}{0pt}{1.8ex plus 1ex minus .2ex}{0.8ex plus .2ex}

% Leave Quarto's own title / author / affiliation block untouched (it renders
% multiple authors as an internal array, which a titling override would break).
% The document title still reads in Georgia + teal because the main font and the
% heading styles below carry the house look.

% --- Links in terracotta (accent), no boxes ---
\usepackage{hyperref}
\hypersetup{colorlinks=true, allcolors=houseTerracotta,
            linkcolor=houseTerracotta, citecolor=houseTerracotta,
            urlcolor=houseTerracotta}

% --- Body colour stays deep teal (readable on white), captions tracked ---
\AtBeginDocument{\color{houseTeal}}

% --- Figure captions: bold "Figure N." reads slightly toward the accent ---
\usepackage{caption}
\captionsetup{font={small,color=houseTeal!85!black},
              labelfont={bf,color=houseTerracotta}, labelsep=period,
              justification=raggedright, singlelinecheck=false}

% Keep paragraphs left-aligned (never justified-centered), generous leading.
\usepackage{ragged2e}
\setlength{\parskip}{0.6em}
\setlength{\parindent}{0pt}
\linespread{1.12}
